\chapter{Introduction}
\label{ch:introduction}

Music plays a fundamental role in human culture, and with the digitization of music, there has been an increasing need for automated systems that can analyze, classify, and organize music content. The ability to automatically recognize musical instruments in audio recordings is of significant value in various applications. For instance, it can enable enhanced music browsing experiences, where users can search and organize music based on the instruments they contain. This metadata can also be used to improve music education tools, assist in music production, and contribute to research in musicology. Moreover, instrument recognition can be a critical component in more complex systems such as automatic music transcription and source separation.

Recognizing musical instruments in audio signals is a challenging task due to the complex and highly variable nature of musical sounds. Instruments can produce sounds with a wide range of timbres and pitches, and their sonic characteristics can change based on playing techniques and recording conditions.

A central aspect of music instrument recognition is the choice of feature representations used to describe audio signals. These features should capture the relevant information in the audio signal that allows different instruments to be distinguished from each other.

One approach to feature extraction involves designing features based on our understanding of music and audio signals. A notable example of this approach is the Scattering Transform, a time-frequency representation that extends wavelet transforms to capture higher-order modulation spectra in signals. Introduced by Joakim And\'{e}n and St\'{e}phane Mallat~\cite{anden2014deep}, the Scattering Transform has shown promise in various audio signal processing tasks due to its ability to capture translation-invariant and stable representations of signals.

An alternative approach to feature extraction is to learn features directly from data using machine learning techniques. One such method is MusicNN, a set of pre-trained deep convolutional neural networks for music tagging, which can be used to extract features from audio signals~\cite{pons2019musicnn}. Specifically, this thesis utilizes the penultimate layer of MusicNN, which provides a rich and discriminative representation of the audio signal that is learned during the training of the neural network.

The primary goal of this thesis is to study and compare the effectiveness of designed features, specifically the Scattering Transform, and learned features, specifically those extracted using MusicNN, in the context of solo instrument recognition. A key question this thesis aims to answer is whether these two types of features contain complementary information. By concatenating these two types of features, we not only explore the potential for improved instrument recognition performance but also investigate whether designed and learned features capture different aspects of the audio signal.

The remainder of this thesis is organized as follows: Chapter~\ref{ch:literature} provides a survey of related work. Chapter~\ref{ch:background} provides the necessary background on signal processing, feature representations, and machine learning methods. Chapter~\ref{ch:method} describes the methods and experimental setup used in this study. Chapter~\ref{ch:results} presents the results of our experiments and provides a detailed analysis. Chapter~\ref{ch:discussion} discusses the implications of our findings, and Chapter~\ref{ch:conclusions} concludes the thesis and suggests directions for future work.
